% =====================================================================
\chapter{和の観測量の分散の導出}
\label{app:sumvar}
% =====================================================================

命題\ref{prop:sumvar}を導く。観測量を $O=\sum_kc_kP_k$、台を $A_k$、各項の文字を $a^{(k)}_q$($q\in A_k$)とする。

\section{1 つの設定での推定値}
設定 $u$ で $n_m$ 回測り、$t$ 番目のショットの結果を $s^{(t)}_q$ とする。項 $k$ の推定値は
\begin{equation}
  X_k(u)=3^{\lvert A_k\rvert}\,\mathbf 1_k(u)\,\bar m_k,\qquad
  \mathbf 1_k(u)=\mathbf 1[\forall q\in A_k:\ b_q=a^{(k)}_q],\qquad
  \bar m_k=\frac1{n_m}\sum_{t=1}^{n_m}\prod_{q\in A_k}s^{(t)}_q
\end{equation}
で、観測量の推定値は $X(u)=\sum_kc_kX_k(u)$ である。CRM では $X^{\rm CRM}(u)=\sum_kc_k3^{\lvert A_k\rvert}\mathbf 1_k(u)(\bar m_k-y_k)+\sum_kc_ky_k$ となる。

\section{2 次のモーメント}
$\mathbb E[X(u)^2]=\sum_{k,l}c_kc_l\,3^{\lvert A_k\rvert+\lvert A_l\rvert}\,\mathbb E\bigl[\mathbf 1_k\mathbf 1_l\,\bar m_k\bar m_l\bigr]$ を計算する。
\begin{itemize}
  \item \textbf{基底の一致.} $\mathbf 1_k\mathbf 1_l=1$ となるのは、$A_k\cup A_l$ のすべての量子ビットで基底が指定どおりのときである。$A_k\cap A_l$ の上で $P_k$ と $P_l$ の文字が違えば(両立しなければ)ありえない。両立するとき、その確率は $3^{-\lvert A_k\cup A_l\rvert}$ である。$\lvert A_k\rvert+\lvert A_l\rvert-\lvert A_k\cup A_l\rvert=\lvert A_k\cap A_l\rvert$ なので、前の因子と合わせて $3^{\lvert A_k\cap A_l\rvert}$ が残る。
  \item \textbf{ショットの平均.} 基底が両方に一致した設定では、1 ショットの $\prod_{q\in A_k}s_q$ と $\prod_{q\in A_l}s_q$ は $P_k$ と $P_l$ の同時固有値で、その積は $P_kP_l$ の固有値である(両立する項は可換で、積は共通の量子ビットで恒等になる)。同じショットどうしの積の平均は $\ev{P_kP_l}$、別々のショットどうしは独立で $x_kx_l$ なので
  \begin{equation}
    \mathbb E[\bar m_k\bar m_l\mid\text{一致}]=\frac1{n_m}\ev{P_kP_l}+\Bigl(1-\frac1{n_m}\Bigr)x_kx_l.
  \end{equation}
\end{itemize}
したがって
\begin{equation}
  \mathbb E[X(u)^2]=\sum_{k,l\ \text{両立}}c_kc_l\,3^{\lvert A_k\cap A_l\rvert}\Bigl[\frac{\ev{P_kP_l}}{n_m}+\Bigl(1-\frac1{n_m}\Bigr)x_kx_l\Bigr]
\end{equation}
で、平均 $\mathbb E[X(u)]=\sum_kc_kx_k$ の 2 乗を引くと式\eqref{eq:sumvar-std}になる。

\section{CRM}
CRM では $\bar m_k$ を $\bar m_k-y_k$ に置き換える($y_k$ は定数)。一致した設定で
\begin{equation}
  \mathbb E[(\bar m_k-y_k)(\bar m_l-y_l)\mid\text{一致}]=\frac{\ev{P_kP_l}-x_kx_l}{n_m}+\Delta_k\Delta_l
\end{equation}
となる(ショットの共分散 $(\ev{P_kP_l}-x_kx_l)/n_m$ と、平均のずれ $\Delta_k\Delta_l$ の和)。定数項 $\sum_kc_ky_k$ は分散に寄与しないので、平均 $\sum_kc_k\Delta_k$ の 2 乗を引いて式\eqref{eq:sumvar-crm}を得る。

単一の項 $P$ では $\ev{P^2}=1$、$3^{\lvert A\cap A\rvert}=3^{\lvert A\rvert}$ なので、式\eqref{eq:var-std}・\eqref{eq:var-crm}に戻る。
